\documentclass[10pt,a4paper]{article}
\usepackage[left=2.5cm,right=2.5cm,top=2.5cm,bottom=2.5cm,]{geometry}
\usepackage{amsmath,amssymb,amsthm}
\usepackage{mathrsfs}
\newtheorem{theorem}{Theorem}[section]
\newtheorem{lemma}[theorem]{Lemma}
\newtheorem{proposition}[theorem]{Proposition}
\usepackage{setspace}
\setlength{\parindent}{0pt}
\onehalfspacing
\begin{document}
I will share some interesting problem here.\\
Q: construct a set $X$ equal $\mathcal{O}(G\curvearrowright X)=|G|$ and the group action is faithful.\\
A: I will present a method based on coset decomposition. For $H\le G$, the left coset space is $G/H=\{xH\mid x\in G\}$.
Consider act on $G$, so $(\forall g,x\in G)g\cdot(xH)=(gx)H$. Obviously, $G/H$ is an obrit and $|G/H|=[G:H]$.
The kernel of this action is $\ker(G\curvearrowright G/H)=\bigcap\limits_{g\in G}gHg^{-1}=:\operatorname{Core}_G(H)$.
The essence of this problem lies in the gluing of coset space. Specificly, let $\{H_i\}_{i\in I}$, $X=\bigsqcup\limits_{i\in I}G/H_i$.
Every $G/H_i$ is an obrit. we only need $\ker(G\curvearrowright G/H_i)=\{1\}\land|I|=|G|$. A key observation is $(\exists H_i)\operatorname{Core}(H_i)=\{1\}$ not $\forall H_i$.
\\
A important truth is $\mathcal{O}=|G|=|X|\iff$ every points are fixed points $\iff$ trivial action.
In the other word, if you want the action is faithful, then $|X|>|G|$. What is this lower bound? if we split $X$ with orbit, 
then $X\cong\bigsqcup\limits_{i=1}^{|G|}G/H_i$, in order to ensure faithful $\bigcap\limits_{i=1}^{|G|}\operatorname{Core}(H_i)=\{1\}$.
so $|X|=\sum[G:H_i]\ge1$. Faithfulness only requires ensuring that one orbit is non-trivial. $|X|\ge(|G|-1)\times1+2=|G|+1$. so $|G|=2$. Interestingly,
Only $\mathbb{Z}/2\mathbb{Z}$ can equal.
\\
\\
Q: Does there exist a function $f(x)$ defined on $\mathbb{R}/\{0\}$ such that $f(f(x)) = x^{-1}$? If so, provide a construction for $f$.\\
A: a beautiful construct: $f(x)=\begin{cases}-x&x>0\\-\frac{1}{x}&x<0\end{cases}$. orz
\\
\\
Q: $f$ is continuous in $[0,1]$ and $g$ is $g$ is an integrable function on $\mathbb{R}$ period $1$. Proof \[\lim\limits_{n\to+\infty}\int_{0}^{1}f(x)g(nx)dx=(\int_{0}^{1}f(x)dx)(\int_{0}^{1}g(x)dx)\]
\begin{proof}
A: Consider split the interval $[0,1]$ into $n$-parts, so $[0,1]=\bigcup\limits_{k=1}^{n}[\frac{k-1}{n},\frac{k}{n}]$.
\[\int_{0}^{1}f(x)g(nx)dx=\sum\limits_{k=1}^{n}\int_{\frac{k-1}{n}}^{\frac{k}{n}}f(x)g(nx)dx\]let $t=nx$, $dx=\frac{1}{n}dt$.
$\forall\varepsilon>0,(\exists\delta>0)|x_1-x_2|<\delta\implies|f(x_1)-f(x_2)|<\varepsilon$. $\forall x\in[\frac{k-1}{n},\frac{k}{n}],|x-\frac{k}{n}|\le\frac{1}{n}<\delta$.
\[|f(x)-f(\frac{k}{n})|<\varepsilon\]
\[\Delta x=|\int_{0}^{1}f(x)g(nx)dx-\sum\limits_{k=1}^{n}f(\frac{k}{n})\int_{\frac{k-1}{n}}^{\frac{k}{n}}g(nx)dx|\]
\[\Delta x=|\sum\limits_{k=1}^{n}\int_{\frac{k-1}{n}}^{\frac{k}{n}}(f(x)-f(\frac{k}{n}))g(nx)dx|\]
\[\Delta x<\varepsilon\int_{0}^{1}|g(nx)|dx\]
\[\varepsilon\int_{0}^{1}|g(nx)|dx=\frac{1}{n}\varepsilon\int_{0}^{n}g(t)dt\]
$n\to+\infty$, $\Delta x\to0$
\[\lim\limits_{n\to\infty}\int_{0}^{1}f(x)g(nx)dx=\lim\limits_{n\to\infty}\sum\limits_{k=1}^{n}f(\frac{k}{n})\int_{\frac{k-1}{n}}^{\frac{k}{n}}g(nx)dx\]
\[(\lim\limits_{n\to\infty}\sum\limits_{k=1}^{n}f(\frac{k}{n}))(\int_{0}^{1}g(nx)dx)=\lim\limits_{n\to\infty}\sum\limits_{k=1}^{n}f(\frac{k}{n})\frac{1}{n}(\int_{0}^{1}g(nx)dx)=(\int_{0}^{1}f(x)dx)(\int_{0}^{1}g(x)dx)\]
\end{proof}
Q:what is the finite-set? what is the infinite-set?\\
A:Formlly, a set $S$ is finite if and only if there exists a map(bejective) $\varphi:S\to\{1,2,3,\cdots,n\}$.for infinite, it's opposite.
\\
\\
Q:Consider a set $S$, if it's a finite-set, it can be defined? if it's a infinite-set, it cannot be defined?
A:the domain of universe is the set theory universe $\mathscr{U}$, a set $A$, $A\subseteq M$ can be defined if and only if there exists a function $\pi(x,y_1,y_2,y_3,\cdots)\rightarrow A=\{y\in M\mid\mathscr{U}|=\pi(y,\bar{x})\}$.
There is no nescessary connection between finiteness and definability. There exists set which is finite but cannot be defined. such $\langle\pi\rangle$ in field $\langle\mathbb{R},+,\cdot,0,1\rangle$. Also, there exists definability but infinite-set. such both even number in $\mathbb{N}$.
Specificly, for a finite set, it can defined with parameters. for example, let $A=\{a_1,a_2,\cdots,a_n\}$,the formula $y=a_1\lor y=a_2\lor y=a_3\cdots\lor y=a_n$.
\\
\\
Q:show that $H$ is the subgroup of $G$ and  $\exists\pi:G\to H$, $\pi\mid_H=\mathrm{id_H.}$, then $G=\ker\pi\rtimes H$.
\begin{proof}
   $\forall h(h\in H\rightarrow \pi(h)=h)$. $\ker\pi\cap H=\{1\}$, $\ker\varphi\trianglelefteq G$. $\forall g\in G$, $g=(g\cdot(\pi(g))^{-1})\cdot\pi(g)$.
    $G=\ker\pi\rtimes H$. more generally, we have 
    \[1\longrightarrow\ker\varphi\overset{\iota}\longrightarrow\ G\longrightarrow H\overset{\pi}\longrightarrow1\]
\end{proof}
Q: Given group $G=\prod\limits_{p\in P}\mathbb{Z}/p\mathbb{Z}$, $T=\bigoplus\limits_{p\in P}\mathbb{Z}/p\mathbb{Z}$, show that $G\cong G/T\oplus T$ is false. 
\end{document}

